\documentclass[12pt, a4paper]{article}

\usepackage{graphicx}

\textheight=210mm
\textwidth=155mm

\hoffset=-1mm
\voffset=-10mm

%% for corrections 
\usepackage{color}
\newcommand\bb{\begin{color}{blue}}
\newcommand\br{\begin{color}{red}}
\newcommand\eb{\end{color}\ }

\begin{document}

\begin{center}
{\bf Response to Referee~\#3's comments}\\
on paper ``Power spectrum scaling as a measure of critical slowing down''\\[5pt]
by {\it J.~Prettyman}
\end{center}

\vspace{3mm}

\noindent We would like to thank the Referee for their time in reviewing this manuscript and for the comments which have informed the revisions made. 

An application of the method is found in [Prettyman et al. 2019]\footnote{Prettyman, J., Kuna, T., and Livina, V. Generalized early warning signals in multivariate and gridded data with an application to tropical cyclones. Chaos: An Interdisciplinary Journal of Nonlinear Science 29, 7 (2019), 073105.}. During this revision we attempted to find another suitable application, while time constraints meant that we could not conduct new research to identify tipping points in different systems. We have therefore included a brief “application” section (Sec. 6) where we apply the PS indicator method to paleo temperature proxy records from GISP2, previously subject to tipping point analysis by [Lenton et al. 2012]\footnote{‌Lenton, T.M., Livina, V.N., Dakos, V. and Scheffer, M., 2012. Climate bifurcation during the last deglaciation?. Climate of the Past, 8(4), pp.1127-1139.} using other methods. The tipping point in these data occur over a period of about 2000 years, or about 100 data points. This is not good enough that the PS indicator is able to provide an early warning signal, but we note that the indicator does increase as the tipping point occurs (albeit not before) which further demonstrates the relationship to the indicators used in the above-mentioned paper (namely DFA and ACF-1), and provides evidence for the presence of slowing-down in the data.

It is our hope that publication of this paper will allow the use of this novel tipping point indicator, possibly alongside other methods, in a greater variety of Earth system applications. We suggest at the end of Sec. 5.2 that modern meteorological records, which often provide hourly data, may provide a suitable application for the PS indicator in climate-linked tipping events over recent decades, such as coral bleaching. This will be an interesting direction for future research but was not possible to undertake in the time spent revising this manuscript. 

\begin{itemize}

\item[\bf Comment 1.]  %%% DONE
{\sl It would be great to embed the results, discussion and implications into the context of real-world systems, to better inform when/how to apply this method to systems of interest.
}

\item[\bf Response.] We provide a brief application section (Sec. 6) in which the PS indicator is applied to a previously-studied paleo-temperature record. The indicator does not provide an early warning signal, although it does provide evidence for the presence of slowing-down in the data. We suggest at the end of Sec. 5.2 that modern meteorological records, which often provide hourly data, may provide a suitable application for the PS indicator in climate-linked tipping events over recent decades, such as coral bleaching.

\item[\bf Comment 2.]
{\sl The abstract could be clearer and more friendly to non-mathematical readers.
}

\item[\bf Response.] We have tried as well as possible to demystify the abstract whilst retaining the meaning. In particular we have revised the line ``considering the CSD phenomenon directly''

\item[\bf Comment 3.] %%% DONE
{\sl The derivation of Eq. 14 seems a bit too oversimplified.
}

\item[\bf Response.] We agree that the derivation of equation 15 was confusing, we hope that the wording, which clarifies the origin of Eqn. 14 (as the Fourier transform of the Brownian process $W(t)$), is now easier to follow. 

%%%%%%%%%%%%%%%%%%%%%%%%%%%%%%%%%%%%%%%%%%%%%%%%%%%%%%%%%%%%%%%%%%

\end{itemize}

\vspace{7mm}

\noindent
We hope that the revised manuscript is now suitable for publication
in the ERL special issue.

\vspace{5mm}

\noindent
Yours sincerely,

\vspace{1mm}

\noindent
T.~Kuna, V.~N.~Livina, and J.~Prettyman

\end{document}

